\begin{abstract}
\noindent
Serving large language models at scale is a distributed systems problem
whose overload behavior is governed at the ingress. We conduct an
adapted-PRISMA systematic review of fifty-five admission and
admission-adjacent works published between 2022 and 2026 across OSDI,
SOSP, NSDI, ASPLOS, EuroSys, MLSys, SoCC, ICML, NeurIPS, HPCA and
arXiv, and report three findings. First, admission sits at a different
layer from the one the serving literature studies: of nine LLM serving
papers audited across a seven-venue slice, all place the overload
response inside the engine, through batching, migration, placement,
autoscaling or memory management, and none takes a decision at the
ingress.
Second, of sixteen admission-primary papers, no two share a baseline
set: policies are measured against the absence of admission control and
almost never against one another. Third, placing all sixteen on our
axes succeeds on six of seven properties and fails on one. Unit,
setting, objective, the quantity a policy watches, the metadata it
assumes, and whether it is deployable in front of an unmodified engine
are all recoverable from the published text; the observed range of the
quantity the policy reads is reported by none of them. Selection
therefore fails on a single missing fact, and it is the fact that
determines whether a reported improvement is interpretable at all.
We contribute a taxonomy of unit, setting,
SLO awareness and signal, with a master classification carrying
GRADE-style evidence ratings; formal definitions of the
admission-decision, goodput-under-admission and fairness-index
primitives; seven reporting items whose inconsistency blocks
cross-paper comparison; a mapping of eight public production incidents
to admission failure modes; and seven open problems, each grounded in
published peer review or in a primary source's own limitations.
\end{abstract}
